% Generated from project outputs. Do not edit by hand.
\begin{table}[!htbp]
\centering
\begin{threeparttable}
\caption{Including remittances by subsample: 2021-2024}
\label{tab:descriptive-2021-2024-subsamples-y4}
\small
\setlength{\tabcolsep}{4pt}
\renewcommand{\arraystretch}{1.15}
\begin{tabularx}{\linewidth}{@{}Xrrrrr@{}}
\toprule
Sample & N & Mean t & SD t & Mean t+1 & SD t+1 \\
\midrule
General & 276 & 3,075.5 & 1,623.8 & 3,444.5 & 1,708.2 \\
Men & 276 & 3,445.1 & 1,803.0 & 3,890.1 & 1,855.4 \\
Women & 276 & 2,488.1 & 1,621.5 & 2,789.1 & 1,736.1 \\
Urban & 276 & 3,268.8 & 1,685.2 & 3,615.5 & 1,741.8 \\
Rural & 276 & 2,756.1 & 1,752.1 & 3,148.9 & 2,021.4 \\
Indigenous & 276 & 2,740.5 & 1,896.6 & 3,230.4 & 1,963.2 \\
Non-Indigenous & 184 & 3,206.7 & 1,666.0 & 3,465.4 & 1,648.1 \\
Bottom 99 percent & 276 & 2,882.0 & 1,381.8 & 3,233.7 & 1,431.8 \\
Bottom 90 percent & 276 & 2,393.2 & 956.7 & 2,674.6 & 962.1 \\
Bottom 50 percent & 276 & 1,407.2 & 556.2 & 1,570.8 & 544.2 \\
\bottomrule
\end{tabularx}
\begin{tablenotes}[flushleft]
\footnotesize
\item Monthly income in quetzales. Unweighted descriptive statistics of survey-weighted cohort means. N counts cohort-transition rows, not individuals or independent cohorts. SD is dispersion across cohort means, not the standard error of a mean. Both incomes must be observed. Subsamples overlap; bottom-income groups are cumulative.
\end{tablenotes}
\end{threeparttable}
\end{table}
